\documentclass[11pt]{article}
\usepackage[T1]{fontenc}
\usepackage[utf8]{inputenc}
\usepackage{amsmath,amssymb,amsthm,booktabs}
\usepackage[margin=1in]{geometry}
\usepackage{hyperref,xurl}
\newtheorem{theorem}{Theorem}
\newtheorem{lemma}{Lemma}
\title{Octahedral facet blocks in a specified\linebreak ideal right-angled five-dimensional spectrum}
\author{Coxeter Polytopes Frontier Project}
\date{18 October 2026}
\begin{document}
\maketitle
\begin{abstract}
Restrict the four-facets of an entirely ideal right-angled hyperbolic
five-polytope to the regular24-cell, its facet double, and one explicit
sixteen-cell polytope. The last type splits into six-facet octahedral
blocks under adjacency along24-vertex ridges. A block has at most two
common adjacent facets. A coloring theorem then excludes the previously
surviving count of30 sixteen-cell facets, improving the necessary counts
to36 such facets and84 doubles. The spectrum itself remains unexcluded;
no new dimension ceiling or geometric example is asserted.
\end{abstract}

\section{Scope and exact input}
Let $C$ be the regular ideal right-angled hyperbolic24-cell and $D$ its
reflection double across one facet. Let $T$ be the particular polytope
obtained from $D$ by successive reflection doubles across the retained
facets with native unit normals
\[
 (1,0,1,-1,0),\quad(1,0,1,0,-1),\quad(1,0,1,1,0)
\]
in metric $\operatorname{diag}(-1,1,1,1,1)$. The inherited doubling
argument proves the halfspace-union identity and finite volume at each
step. Its frozen exact certificate lists all198 walls and258 vertices
of $T$. This paper assumes every four-facet of $P$ is isometric to
$C,D$, or this particular $T$. It makes no classification assumption
about other ideal right-angled four-polytopes.

Write $x,y,z$ for these three facet counts. The predecessor carrier and
cusp-count theorems, with the same spectrum hypothesis, give
\begin{equation}\label{eq:counts}
 y=3z-24,\qquad 6\mid z,\qquad z\ge30.
\end{equation}
They are internal proved inputs documented separately, rather than new
external assumptions. Our only additional published theorem is the
Andr\'asfai--Erd\H{o}s--S\'os coloring theorem~\cite{AES}.

Exactly four ridges of $T$ have24 vertices. Their native unit normals,
in saved order5,14,30,37, have Gram matrix
\begin{equation}\label{eq:gram}
 \begin{pmatrix}1&0&0&-1\\0&1&-1&0\\0&-1&1&0\\-1&0&0&1\end{pmatrix}.
\end{equation}
Both opposite-pair sums are positive multiples of the same ideal ray,
saved vertex85. Each of these three-ridges has26 polygon faces:
twelve triangles, twelve quadrilaterals and two hexagons. Precisely
the two hexagons are its intersections with the other24-vertex ridges
of $T$. Complete rational reconstruction of all104 records verifies
these statements. Neither $C$ nor $D$ has a24-vertex ridge.

\section{The octahedral partition}
\begin{lemma}\label{lem:blocks}
The graph $H$ on $T$ facets, with edges corresponding to24-vertex
common ridges, is a disjoint union of induced $K_{2,2,2}$ graphs.
Every block induces the same octahedral graph in the full facet graph.
\end{lemma}
\begin{proof}
Every $T$ has exactly four such ridges, with distinct neighboring
facets, all of type $T$. At a shared24-vertex ridge, the two hexagons
are intrinsically identifiable. Seen from either adjacent $T$, they
are exactly the intersections with its other24-vertex ridges. Thus
the two corresponding neighboring facets are adjacent in $H$ to both
ends of this edge. A two-face is contained in exactly three ambient
facets, since its finite relative interior has a right-angled simplex
link; the correspondence therefore identifies actual facets.
Matrix~\eqref{eq:gram} says the neighborhood of each vertex induces
a four-cycle. The opposite pairs are asymptotic, so cannot be adjacent
even in the full facet graph.

For completeness, a finite four-regular locally four-cycle graph is
octahedral component by component. Give the neighbors of $v$ cyclic
names $a,b,c,d$. The fourth neighbor $w$ of $a$ lies outside this set.
The four-cycle in its neighborhood forces $w$ adjacent to $b,d$ and
not to $v$. Now the neighborhood of $b$ is $v,a,c,w$; its four-cycle
forces $w$ adjacent to $c$. All six degrees are filled by the resulting
octahedron. The same description of opposite pairs inside its facets
excludes extra diagonals in the full graph.
\end{proof}

\begin{lemma}\label{lem:common}
At most two facets are adjacent to every member of an octahedral block.
\end{lemma}
\begin{proof}
Use ambient Lorentz signature$(5,1)$. Choose a block normal $n$ and
its four neighbors $r_i$ in $H$. Their span has two positive directions
and a one-dimensional null radical $v$, by~\eqref{eq:gram} and its
opposite-pair sums. Adjoining $n$ gives a four-dimensional space $U$
with signature$(3,0,1)$. The sixth normal $m$ is orthogonal to the
four $r_i$ and satisfies $\langle n,m\rangle=-1$: the latter follows
from an adjacent facet's opposite pair in~\eqref{eq:gram}.
Hence $m+n\in U^\perp$ is null. This orthogonal complement consists
of one positive direction and radical $v$, so $m+n$ is a multiple of
$v$ and all six block normals span $U$.

A common adjacent facet has unit normal in $U^\perp$. Choose a unit
positive vector $e$ there. Such normals have form $e+tv$ or $-e+tv$.
Two with the same sign have product one, whereas distinct outward
normals of an acute polytope have nonpositive product. At most two
are possible.
\end{proof}

\section{The coloring obstruction}
\begin{lemma}\label{lem:clique}
The full facet-adjacency graph of an entirely ideal right-angled
hyperbolic five-polytope is $K_5$-free.
\end{lemma}
\begin{proof}
Five pairwise adjacent facets would have orthonormal unit normals
$n_1,\ldots,n_5$. For a strict timelike interior vector $q$, put
\[
 p=q-\sum_i\langle q,n_i\rangle n_i.
\]
Then $\langle p,p\rangle=\langle q,q\rangle-
\sum_i\langle q,n_i\rangle^2<0$, and $p$ has the same time orientation
as $q$. For any other outward normal $m$,
\[
 \langle p,m\rangle=\langle q,m\rangle-
 \sum_i\langle q,n_i\rangle\langle n_i,m\rangle
 \le\langle q,m\rangle<0.
\]
Thus $p$ normalizes to an actual finite vertex on the five selected
facets, contrary to ideality.
\end{proof}

\begin{theorem}
Under the stated three-type facet restriction, $z\ge36$ and $y\ge84$.
\end{theorem}
\begin{proof}
Let $G$ be the full ridge-adjacency graph on the $z$ facets of type $T$.
Every $T$ has22 ridges with14 or24 vertices, each shared with a distinct
$T$, so $\delta(G)\ge22$. Suppose $z=30$.
By Lemma~\ref{lem:clique}, $G$ is $K_5$-free. Theorem1.1 of~\cite{AES},
with its parameter $r=5$, gives four-colorability when
$\delta(G)>8z/11$, which holds since $22>240/11$.

Each nonempty color class has size at most eight, since its vertices
need22 neighbors outside it. Hence four nonempty classes have sizes
$(6,8,8,8)$ or$(7,7,8,8)$. Vertices in an eight-element class are
adjacent to all vertices outside their class. Vertices in a seven-element
class miss at most one vertex outside their class.

Consider one of the five disjoint octahedral blocks from
Lemma~\ref{lem:blocks}. Three colors would make its three opposite
pairs monochromatic. If the omitted class has size eight, all eight
vertices are common neighbors. If the sizes are$(6,8,8,8)$ and the
omitted class has size six, all six are common neighbors. In the
remaining case the omitted class has size seven and only the two
block vertices in the other seven-element class can miss its members,
at most one each. There are at least five common neighbors. Each
case contradicts Lemma~\ref{lem:common}.

Thus the block uses four colors, with exactly one bichromatic opposite
pair. Neither endpoint of that nonedge can be in an eight-element
class. The part sizes must therefore be$(7,7,8,8)$, with the split
pair in the two seven-element classes. The other pairs use the two
eight-element classes. Every block has color counts$(1,1,2,2)$.
Five blocks would give$(5,5,10,10)$ instead of$(7,7,8,8)$, a contradiction.

This excludes $z=30$. Equation~\eqref{eq:counts} now gives $z\ge36$
and $y=3z-24\ge84$.
\end{proof}

\section{Reproduction and remaining scope}
The new standard-library exact-fraction producer is
\path{scripts/fresh_ideal_ra5_octahedral_blocks_20261018.py}; replay its
saved certificate with \texttt{--verify}. It reconstructs the four
ridge Gram and polygon-face data and exhausts all4096 four-colorings
of a six-element block for each possible part-size vector. The full
written proof, input hashes and verification scope are recorded in
\path{results/ideal_ra5_octahedral_blocks_20261018.md} and the associated
reproduction map. No optimizer, unbounded search, Lean compilation,
or classification of unlisted facet types is used. The universal
geometric arguments remain written proofs awaiting formalization.

\begin{thebibliography}{1}
\bibitem{AES} B.~Andr\'asfai, P.~Erd\H{o}s and V.~T.~S\'os,
On the connection between chromatic number, maximal clique and minimal
degree of a graph, \emph{Discrete Mathematics}8 (1974),205--218.
\url{https://real.mtak.hu/110445/1/1974-18.pdf}.
\end{thebibliography}
\end{document}
